\begin{statement}{C.2.UN.68}{Proved}{Rational toric component pairs and the exact dual torus}
Let $f:X\to\Delta$ have the full reduced simple-normal-crossing
central fibre $S=\sum_{i=1}^{v}V_i$ of C.2.UN.65, and let its
canonical evaluation orders be equal. Suppose that $e(S)>0$.
Then every $V_i$ is rational and its actual boundary
$D_i=\sum_{j\ne i}V_i\cap V_j$ is the full toric boundary
for an actual algebraic two-torus action on $V_i$.
The dual complex of this given fibre is a closed real
two-torus. If $r$ counts its connected double curves and
$t$ its triple points, then
\[
                  r=3v,\qquad t=2v,\qquad e(S)=2v.
\tag{C.TC.1}
\]
In particular $v\ge3$ and $e(S)\ge6$ in the embedded
simple-normal-crossing setting of the statement.

Together with C.2.UN.65 this identifies every component
pair and the dual topological space of a canonical-flat
torus degeneration. The assertion concerns the actual
given degeneration and its actual component pairs.
No smoothing of arbitrary proposed toric gluing data
is inferred from these necessary comparisons.
\end{statement}

\begin{proof}
\par\medskip\noindent\textbf{A rational boundary cycle forces a rational component.}
Use the actual three-form $\Omega/f^k$ of UN65.
It gives
\[
               K_{V_i}+D_i=0,\qquad
               K_C=\OO_C(-T_C)
\tag{C.TC.2}
\]
by its component residue and the principal-fibre line maps.
Thus a boundary component is either elliptic and disjoint
from the remaining boundary, or rational and meets it
at exactly two transverse points. Consequently the
rational boundary components form disjoint nodal cycles.
The positive Euler identity $e(S)=t$ of UN65 supplies
at least one such cycle.

We check how an actual point blowdown on a component
acts on this divisor. Let $E$ be a smooth rational
$(-1)$-curve and $\pi:V\to V'$ its ordinary blowdown.
Adjunction gives
\[
                         D\cdot E=1.
\tag{C.TC.3}
\]
If $E$ is not in $D$, it meets exactly one smooth
boundary point transversely. It cannot meet a node,
which would contribute at least two to (C.TC.3).
The image $D'=\pi_*D$ is again reduced nodal and
$K_{V'}+D'=0$. Locally this is the ordinary blowup
at a smooth point of $D'$ with $D$ its strict transform.

If $E$ is in $D$, then
$(D-E)\cdot E=2$. Its two intersection points are
distinct and transverse. The blowdown identifies
their two transverse tangent directions and creates
one ordinary node on $D'$. The two branches may
belong to the same irreducible curve. This is the
actual inverse of the blowup of a boundary node;
again $K_{V'}+D'=0$ by the ordinary canonical
Jacobian formula. In particular a rational cycle
persists. When it has become an irreducible nodal
rational curve, it cannot be the smooth exceptional
curve of a further ordinary point blowdown.

These two maps preserve the separation between
the rational cycles and the elliptic boundary
components. Contracting the point blowups to
a minimal model therefore retains a rational
cycle whenever one was present. Each component
has $\kappa=-\infty$, because $-K=D$ is
nonempty effective. A nonrational ruled minimal
model has rational curves only in its ruling
fibres: a map from $\PP^1$ to the positive-genus
base is constant. Those fibres are mutually
disjoint smooth rational curves and cannot
contain a nodal rational cycle. The other
nonrational possibility is the Hopf minimal
model of the actual class-VII component,
by UN19. A rational curve on a Hopf surface
would lift to $\C^2\setminus\{0\}$, where
its coordinate functions are constant. It
has no rational cycle either. Hence a
component whose boundary contains a rational
cycle is rational.

On a rational component the exact sequence
\[
 0\longrightarrow K_V\longrightarrow\OO_V
                 \longrightarrow\OO_D\longrightarrow0
\tag{C.TC.4}
\]
has $H^0(K_V)=H^1(K_V)=0$, the second vanishing
being Serre dual to $H^1(\OO_V)=0$.
It follows that $H^0(\OO_D)=\C$, so $D$ is
connected. If $D$ contains a rational curve,
it is therefore exactly one rational nodal
cycle. A neighbouring component contains
that same rational double curve and hence
is rational by the preceding argument.
Propagation through the connected fibre
shows that every $V_i$ is rational and
every $D_i$ is one rational cycle.

\par\medskip\noindent\textbf{The actual dual complex and its cohomology.}
Let $\mathcal B$ be the dual complex with
one vertex per $V_i$, one edge per connected
double curve and one triangle per triple
point. Multiple edges and triangles are
retained. Every edge belongs to two
triangles by (C.TC.2), and the link of
each vertex is exactly the boundary cycle
on that component. Thus the actual
topological realization of $\mathcal B$
is a connected closed real surface.

All positive-degree structure-sheaf
cohomology of the rational components
and rational double curves vanishes.
The SNC restriction sequence therefore
identifies, by its alternating restriction
maps, the full complex computing
$H^*(S,\OO_S)$ with
\[
           \C^{v}\longrightarrow\C^{r}
                                  \longrightarrow\C^{t},
\tag{C.TC.5}
\]
the cellular cochain complex of
$\mathcal B$. The principal-fibre
adjunction and the actual volume form
give $\omega_S\simeq\OO_S$.
Serre duality on the compact Gorenstein
surface $S$ consequently gives
\[
             H^2(S,\OO_S)\simeq
                         H^0(S,\OO_S)^*\simeq\C.
\tag{C.TC.6}
\]
Here $H^0(S,\OO_S)=\C$ follows from
compactness, connectedness and reducedness.
Equations (C.TC.5)--(C.TC.6) make
$\mathcal B$ orientable. Proper flatness
gives $\chi(\OO_S)=0$, hence
$\chi(\mathcal B)=0$. A connected closed
oriented surface with this Euler
characteristic is a real two-torus.

The same cochain counts give
$v-r+t=0$. Together with the exact
incidence count $2r=3t$ of UN65,
this proves (C.TC.1). A triple point
in a simple-NC divisor lies on three
distinct embedded components, so
$t>0$ implies $v\ge3$.

\par\medskip\noindent\textbf{An elementary charge test with its actual equality morphisms.}
For a smooth rational surface with an
anticanonical nodal rational cycle
$D=D_1+\cdots+D_n$, allowing an
irreducible nodal rational curve when
$n=1$, put
\[
                        Q(V,D)=e(V)-n=e(V\setminus D).
\tag{C.TC.7}
\]
Indeed a rational nodal cycle has
Euler characteristic $n$, including
the one-component case. We prove
\[
 Q(V,D)\ge0,\qquad Q(V,D)=0
       \quad\Longleftrightarrow\quad
       (V,D)\text{ is an actual toric pair}.
\tag{C.TC.8}
\]

Follow any sequence of ordinary
$(-1)$-curve contractions to a
minimal rational surface. The two
maps already proved after (C.TC.3)
keep an anticanonical nodal cycle.
At a contraction outside the boundary,
$e(V)$ drops by one and $n$ is
unchanged. At a boundary contraction,
both drop by one. Thus
\[
 Q(V,D)=Q(V_{\min},D_{\min})+a,
\tag{C.TC.9}
\]
where $a$ is the number of actual
contracted exceptional curves outside
the current boundary.

On $\PP^2$ the boundary has total
degree three, so $n\le3=e(\PP^2)$.
Equality means three distinct lines;
their nodal cycle condition makes
them nonconcurrent. The invertible
matrix of their linear equations is
an actual projective transformation
to the coordinate triangle.

On a minimal rational ruled surface
$q:\mathbb F_b\to\PP^1$, the total
horizontal degree of $D$ is two.
There are at most two horizontal
components. Each vertical prime is
a complete ruling fibre and meets
the horizontal part in two points
counted with multiplicity. The cycle
has only two incidences at each
horizontal component, so there are
at most two vertical components.
Hence $n\le4=e(\mathbb F_b)$.
Equality forces two degree-one
sections and two ruling fibres,
with the sections disjoint.
Their two quotient line maps split
the ruling bundle. A base coordinate
puts the two marked fibres at zero
and infinity. The resulting actual
pair is
\[
 \left(\PP_{\PP^1}(\OO\oplus\OO(b)),
        \text{two sections plus the fibres over }0,\infty\right).
\tag{C.TC.10}
\]
In a line coordinate the transitions
are $w_\infty=z^{-b}w_0$.
The maps $(z,w_0)\mapsto(\lambda z,\mu w_0)$
and $(z^{-1},w_\infty)\mapsto
(\lambda^{-1}z^{-1},\lambda^{-b}\mu w_\infty)$
give its actual two-torus action
and precisely the four specified
boundary components.

The two minimal computations prove
$Q(V_{\min},D_{\min})\ge0$.
If the initial charge is zero,
(C.TC.9) says that $a=0$ and
the minimal pair is one of the
two explicit toric pairs just
identified. Reversing the recorded
contractions therefore blows up
only boundary nodes. Each such
node is fixed by the torus action;
the invariant ideal $(x,y)$ gives
the actual lifted action on its
blowup charts $y=xu$ and $x=yv$.
The new full boundary is the
toric boundary. Induction gives
an actual torus action on the
original marked pair $(V,D)$.
Conversely a smooth complete
toric surface has one fixed
point per boundary node; its
orbit decomposition gives
$e(V)=n$, proving the reverse
implication in (C.TC.8).

This direct argument agrees with
the actual toric-model comparison
of \cite[Definition 1.2 and Proposition 1.3]{GHK2015},
whose hypotheses are exactly a
smooth rational projective surface
and an anticanonical nodal rational
cycle. The displayed contractions
prove the equality case without
discarding the original pair.

\par\medskip\noindent\textbf{The global charge is zero.}
Write $n_i$ for the number of
curves in $D_i$. Every double
curve belongs to two component
boundaries, so
\[
                         \sum_i n_i=2r=3t.
\tag{C.TC.11}
\]
The exact HRR and boundary-square
calculation in UN65 gives
\[
            \sum_i e(V_i)
             =12(r-t)-3t=3t.
\tag{C.TC.12}
\]
Subtracting (C.TC.11) from
(C.TC.12) yields
$\sum_i Q(V_i,D_i)=0$.
Every summand is nonnegative by
(C.TC.8), so each is zero.
The actual equality morphisms
above make every specified
$(V_i,D_i)$ a toric pair.
This proves the asserted
classification on the given
degeneration.
\end{proof}


% C.2.UN.44 from part_C2_fundamental_groups.tex

